%!TEX root = ./main.tex

\usetikzlibrary{positioning,fit,backgrounds}

\definecolor{ranblue}{RGB}{46,91,186}
\definecolor{coreorange}{RGB}{222,122,36}
\definecolor{ranpurple}{RGB}{109,53,151}
\definecolor{softgreen}{RGB}{73,139,74}
\definecolor{softgray}{RGB}{92,98,105}

\newcommand{\source}[2]{%
  \par\vspace{0.25em}%
  {\tiny\textcolor{gray}{\href{#1}{\underline{#2}}}\par}%
}

% Inline expansion of an acronym at first use: SMS \expand{Short Message Service}
\newcommand{\expand}[1]{\textcolor{softgray}{(#1)}}

% Small gray glossary line, for slides with several acronyms at once.
\newcommand{\glossline}[1]{%
  \par\vspace{0.4em}%
  {\scriptsize\color{softgray}\raggedright#1\par}%
}

\newcommand{\plainterm}[2]{%
  \textbf{#1}\hspace{0.25em}{\small\textcolor{softgray}{#2}}%
}

\newcommand{\ask}[1]{%
  \begin{tcolorbox}[colback=blue!4,colframe=ranblue,boxrule=0.7pt,arc=2pt,
    left=5pt,right=5pt,top=4pt,bottom=4pt]
    \centering\large #1
  \end{tcolorbox}%
}

% Font is deliberately \normalsize, not \large: at \large a one-sentence takeaway
% wraps to two lines and the box grows into whatever is below it.
\newcommand{\takeaway}[1]{%
  \begin{tcolorbox}[colback=orange!7,colframe=coreorange,boxrule=0.7pt,arc=2pt,
    left=5pt,right=5pt,top=4pt,bottom=4pt]
    \centering\normalsize #1
  \end{tcolorbox}%
}

\tikzset{
  netnode/.style={draw=ranblue,fill=blue!5,rounded corners=3pt,
    minimum width=1.9cm,minimum height=0.82cm,align=center,font=\bfseries},
  corenode/.style={draw=coreorange,fill=orange!7,rounded corners=3pt,
    minimum width=1.9cm,minimum height=0.82cm,align=center,font=\bfseries},
  oranode/.style={draw=ranpurple,fill=purple!6,rounded corners=3pt,
    minimum width=1.9cm,minimum height=0.82cm,align=center,font=\bfseries},
  flow/.style={-{Latex[length=2.4mm]},thick,draw=softgray}
}

% Full-width question/answer block: spans the whole frame, no box.
% Use below \end{columns} when the question is meant for the whole slide.
\newcommand{\qa}[2]{%
  \par\vspace{0.55em}%
  {\large\raggedright
    \noindent\hangindent=1.25em\hangafter=1
    \textcolor{ranblue}{\textbf{Q:}}\hspace{0.35em}\textbf{#1}\par
    \vspace{0.35em}
    \noindent\hangindent=1.25em\hangafter=1
    \textcolor{coreorange}{\textbf{A:}}\hspace{0.35em}#2\par}%
}

% Repeating template for "this box is gone in the next diagram" frames.
% The deck's spine: a term never vanishes between two architecture diagrams
% without a frame explaining the pressure that removed it.
% #1 what it did, #2 what broke, #3 what replaced it
\newcommand{\gonebox}[3]{%
  \vspace{0.5em}%
  \begin{columns}[T,onlytextwidth]
    \column{0.30\textwidth}
      {\small\textbf{\textcolor{softgray}{What it did}}\par}
      \vspace{0.4em}
      {\small #1\par}
    \column{0.34\textwidth}
      {\small\textbf{\textcolor{coreorange}{What broke}}\par}
      \vspace{0.4em}
      {\small #2\par}
    \column{0.30\textwidth}
      {\small\textbf{\textcolor{softgreen}{What replaced it}}\par}
      \vspace{0.4em}
      {\small #3\par}
  \end{columns}%
}

% ---------------------------------------------------------------------------
% Tour map for the second half of the deck (added 2026-09-09).
% The lecture's back half is five papers; students felt the seams between them.
% \tourmap{n} draws the five stops with stop n highlighted. Shown once in full
% ("Five places the curve must survive") and reprised at every later stop, so
% no topic ever appears without a visible place on the route.
\tikzset{
  stopbox/.style={draw=softgray!75,fill=white,rounded corners=3pt,
    minimum width=1.72cm,minimum height=0.86cm,align=center,font=\scriptsize},
  herebox/.style={draw=coreorange,fill=orange!12,very thick,rounded corners=3pt,
    minimum width=1.72cm,minimum height=0.86cm,align=center,font=\scriptsize}
}
% #1 highlighted stop, #2 stop index, #3 x position, #4 box label, #5 gloss below
% Gotcha encoded here: a \node cannot sit inside \ifnum directly, so the
% conditional picks a style NAME. Labels are single-line on purpose (line breaks
% passed through macro arguments misbehave inside beamer frames); the tiny gloss
% is its own node under the box instead of a second line in it.
\newcommand{\tourstopnode}[5]{%
  \ifnum#1=#2\relax\def\tourstopstyle{herebox}%
  \else\def\tourstopstyle{stopbox}\fi
  \node[\tourstopstyle] (s#2) at (#3,0) {\textbf{#4}};
  \node[font=\tiny,text=softgray,anchor=north,yshift=-1pt] at (s#2.south) {#5};}
\newcommand{\tourmap}[1]{%
  \begin{tikzpicture}[scale=0.88,transform shape]
    \tourstopnode{#1}{1}{0}{the grid}{allocation}
    \tourstopnode{#1}{2}{2.42}{an adversary}{DISCO}
    \tourstopnode{#1}{3}{4.84}{a train}{350 km/h}
    \tourstopnode{#1}{4}{7.26}{light}{FSO}
    \tourstopnode{#1}{5}{9.68}{a lab bench}{prototype}
    \foreach \a/\b in {s1/s2,s2/s3,s3/s4,s4/s5} \draw[flow] (\a) -- (\b);
    \node[font=\tiny\bfseries,text=coreorange,anchor=south,yshift=1pt]
      at (s#1.north) {YOU ARE HERE};
  \end{tikzpicture}}
